%% =====================================================================
%%  自生长研究智能体 —— 课堂演示
%%  编译方式：Overleaf 中请将 Compiler 设为 XeLaTeX
%%  （ctex 宏包在 XeLaTeX 下对中文支持最完整）
%% =====================================================================
\documentclass[aspectratio=169,11pt]{beamer}

\usepackage[UTF8]{ctex}
\usepackage{amsmath,amssymb}
\usepackage{booktabs}
\usepackage{array}
\usepackage{tikz}
\usetikzlibrary{arrows.meta,positioning,shapes.geometric}
\usepackage{xcolor}

\usetheme{Boadilla}
\usecolortheme{seahorse}
\setbeamertemplate{navigation symbols}{}
\setbeamertemplate{blocks}[rounded][shadow=false]

%% 自定义强调色
\definecolor{accent}{RGB}{20,90,160}
\definecolor{soft}{RGB}{240,244,249}
\definecolor{warn}{RGB}{170,60,40}
\setbeamercolor{structure}{fg=accent}
\setbeamercolor{block title}{bg=accent!12,fg=accent}
\setbeamercolor{block body}{bg=soft}

%% 紧凑列表
\setbeamertemplate{itemize items}[circle]
\setbeamertemplate{enumerate items}[default]
\setlength{\parskip}{2pt}

\renewcommand{\emph}[1]{\textbf{#1}}   % 中文字型无斜体，按排印惯例改着重
\newcommand{\code}[1]{\texttt{\small #1}}
\newcommand{\key}[1]{\textbf{\color{accent}#1}}

\title[自生长研究智能体]{面向构造性组合下界的自生长研究智能体}
\subtitle{以 Ramsey 数 $R(5,5)$ 为检验场景}
\author[曹成博]{曹成博 \quad 学号 2452035}
\institute{人工智能选修课\\课程项目报告}
\date{\today}

\begin{document}

%% ---------------------------------------------------------------- 封面
\begin{frame}
  \titlepage
\end{frame}

%% ---------------------------------------------------------------- 核心指标
\begin{frame}{核心指标}
  \begin{columns}[T,onlytextwidth]
    \begin{column}{0.32\textwidth}
      \begin{block}{86}
        \centering
        单次自主研究运行中的\\工具调用次数
      \end{block}
    \end{column}
    \begin{column}{0.32\textwidth}
      \begin{block}{76.9\%\,$\rightarrow$\,100\%}
        \centering
        留出基准通过率\\（自建工具引入前后）
      \end{block}
    \end{column}
    \begin{column}{0.32\textwidth}
      \begin{block}{46}
        \centering
        覆盖全部模块的\\自动化测试用例数
      \end{block}
    \end{column}
  \end{columns}

  \vspace{0.6em}
  \begin{block}{本报告的主结论}
    所构建的智能体具备\textbf{能力自扩展}机制：其在运行中产生的代码经验，
    经独立校验后可固化为后续运行可直接调用的原生工具；
    该能力增长以留出基准上的通过率曲线度量，而非以单一最优值表征。

    \vspace{0.3em}
    \small\color{black!65}
    数学问题在本项目中承担检验场景的角色，其本身不是研究目的。
  \end{block}
\end{frame}

%% ---------------------------------------------------------------- 问题定位
\begin{frame}{问题定位：智能体与求解器的区别}
  \begin{center}
  \small
  \begin{tabular}{@{}p{0.24\textwidth}p{0.30\textwidth}p{0.32\textwidth}@{}}
    \toprule
    \textbf{维度} & \textbf{一次性求解器} & \textbf{本项目的智能体} \\
    \midrule
    优化对象 & 解空间中的点 & 智能体自身的能力集 \\
    运行间状态 & 每次从零开始 & 继承既有工具与测量经验 \\
    能力表征 & 单一最优目标值 & 随运行次数变化的曲线 \\
    失败处理 & 丢弃本次结果 & 以类型化记录留存，避免重复 \\
    \bottomrule
  \end{tabular}
  \end{center}

  \vspace{0.8em}
  \begin{alertblock}{判据}
    若智能体在运行结束后无法保留新增能力，则其行为等价于一次性求解器。
    \textbf{能力能否跨运行累积，是区分二者的可操作判据。}
  \end{alertblock}
\end{frame}

%% ---------------------------------------------------------------- 为何选 Ramsey
\begin{frame}{检验场景的选择依据}
  \textbf{选择标准}：该领域须提供\textbf{可自动判定的真值}，
  否则无法区分“结果正确”与“智能体认为结果正确”。

  \vspace{0.5em}
  \begin{center}
  \small
  \begin{tabular}{@{}ll@{}}
    \toprule
    任务类型 & 正确性可否自动判定 \\
    \midrule
    文本生成 & 否 \\
    程序合成 & 是（单元测试） \\
    \textbf{给定图的单色团检测} & \textbf{是（有限检查）} \\
    \bottomrule
  \end{tabular}
  \end{center}

  \vspace{0.6em}
  \begin{block}{形式化}
    给定 $n$ 顶点完全图的红蓝二染色，判定其是否含红色 $K_s$ 或蓝色 $K_t$，
    是 $O(n^{s}+n^{t})$ 的有限检查。因此可以为该任务构造
    \textbf{不可争辩的裁决者}，使“正确性”不依赖语言模型的判断。
  \end{block}

  \vspace{0.4em}
  \footnotesize\color{black!70}
  已知界：$43 \le R(5,5) \le 46$。因 43 顶点合法染色存在且删点保持合法性，
  故 $n \le 43$ 全部存在；\textbf{开放靶子仅为 $n=44$ 与 $n=45$}。
\end{frame}

%% ---------------------------------------------------------------- 架构
\begin{frame}{系统架构}
  \begin{center}
  \begin{tikzpicture}[
      box/.style={draw=accent!60,fill=soft,rounded corners=3pt,
                  minimum width=9.6cm,minimum height=0.92cm,
                  align=center,font=\small},
      lbl/.style={font=\scriptsize\color{black!60},align=left},
      arr/.style={-{Stealth[length=2.4mm]},accent!70,thick}
    ]
    \node[box] (rt) {\textbf{Moonshine 运行时}（开源框架，核心逻辑未修改）\\
      \scriptsize 研究循环、类型化研究日志、
      记忆检索、验证门禁、 报告生成};
    \node[box,below=0.68cm of rt] (pack)
      {\textbf{ramsey-pack}：7 个确定性工具、4 份方法论技能、1 个智能体人格};
    \node[box,below=0.68cm of pack] (core)
      {\textbf{marsa} 计算内核、\textbf{growth} 生长回路};

    \draw[arr] (pack) -- (rt);
    \draw[arr] (core) -- (pack);
    \node[lbl,right=0.35cm of rt.east,anchor=north west] {应用层};
    \node[lbl,right=0.35cm of pack.east,anchor=north west] {领域层};
    \node[lbl,right=0.35cm of core.east,anchor=north west] {计算层};
  \end{tikzpicture}
  \end{center}

  \vspace{0.3em}
  \begin{alertblock}{必须补充确定性层的原因}
    该框架原生的验证机制为\textbf{语言模型评审}，其审查对象是证明过程，
    无法判定“给定图不含 $K_5$”这类\textbf{计算性断言}。
    将计算性断言提交语言模型评审，只能得到表述自信但无判别力的回答。
  \end{alertblock}
\end{frame}

%% ---------------------------------------------------------------- 工具层
\begin{frame}{确定性工具层}
  \begin{center}
  \scriptsize
  \begin{tabular}{@{}llp{0.42\textwidth}@{}}
    \toprule
    \textbf{工具} & \textbf{功能} & \textbf{关键性质} \\
    \midrule
    \code{verify\_construction} & 精确验证构造 & 与搜索器\textbf{独立实现}，担任唯一仲裁者 \\
    \code{enumerate\_family}    & 穷举结构族   & \textbf{判定性}：可给出族级排除结论 \\
    \code{search}               & 定预算搜索   & 成功须经独立复核 \\
    \code{violations}           & 统计违反度   & 精确且低开销 \\
    \code{compare}              & 多算子比较   & \textbf{统计纪律自检} \\
    \code{grow}                 & 注册自建工具 & 能力固化 \\
    \code{prior}                & 查询累积经验 & 经验复用 \\
    \bottomrule
  \end{tabular}
  \end{center}

  \vspace{0.5em}
  \begin{columns}[T,onlytextwidth]
    \begin{column}{0.33\textwidth}
      \begin{block}{约束优先于指导}
        \footnotesize
        \code{compare} 直接返回 \code{clear\_winner} 或
        \code{inconclusive}，调用方只能照录。
        诚实性由\textbf{接口返回值}保证，而非由提示词要求。
      \end{block}
    \end{column}
    \begin{column}{0.33\textwidth}
      \begin{block}{验证自动化}
        \footnotesize
        构造合法性是有限检查，交由独立实现判定，
        不以模型输出为依据。
      \end{block}
    \end{column}
    \begin{column}{0.33\textwidth}
      \begin{block}{反馈结构化}
        \footnotesize
        每条结果标注类型（已证明 / 已测量 / 无定论）
        与作用域，三者不得混用。
      \end{block}
    \end{column}
  \end{columns}
\end{frame}

%% ---------------------------------------------------------------- 自主运行实证
\begin{frame}{自主运行的实证结果}
  \textbf{实验设置}：研究模式，迭代预算 10 轮；实际执行
  \key{86 次工具调用}，研究日志 500 行以上。

  \vspace{0.5em}
  \begin{enumerate}
    \item \textbf{工具自造。} 框架的结构族枚举仅覆盖循环图；智能体自行实现
      非循环 Cayley 群枚举器（Gray-code 遍历、分片与断点续跑）。
      首版因逐掩码重建邻接矩阵而超时，其自行定位根因并改为平移构造。

    \item \textbf{自建工具的自校验。} 在使用前设立 5 个锚点
      （$n=38$ 应有解；$n=43,44,45,46$ 应无解），逐一复现引擎判定，
      全部通过后才承认其实现可用。

    \item \textbf{拒绝外推。} 对 23.8\% 的部分扫描，其结论为\emph{“这不是判定，不是排除”}。

    \item \textbf{结论撤回。} 曾给出“强正则图型构造不可能”的论证，
      后经复查发现该论证依赖无效必要条件，遂主动撤回并标注
      \emph{“不得引用为已排除”}。
  \end{enumerate}

  \vspace{0.4em}
  \begin{block}{}
    \small
    上述四类行为共同构成判据：系统具备\textbf{工具构造}、
    \textbf{自我校验}与\textbf{证据不足时拒绝断言}的能力，
    而非仅执行预设流程。
  \end{block}
\end{frame}

%% ---------------------------------------------------------------- 生长回路
\begin{frame}{生长回路：能力跨运行累积的机制}
  \begin{center}
  \begin{tikzpicture}[
      st/.style={draw=accent!55,fill=soft,rounded corners=3pt,
                 minimum width=2.6cm,minimum height=0.95cm,
                 align=center,font=\scriptsize},
      arr/.style={-{Stealth[length=2.2mm]},accent!75,thick},
      note/.style={font=\tiny\color{black!60},align=center}
    ]
    \node[st] (w) {写出工具\\（运行 A）};
    \node[st,right=0.55cm of w] (t) {编写测试};
    \node[st,right=0.55cm of t] (g) {\code{grow}\\\code{register}};
    \node[st,right=0.55cm of g,fill=accent!18] (n) {下次运行\\原生可用};
    \draw[arr] (w) -- (t);
    \draw[arr] (t) -- (g);
    \draw[arr] (g) -- (n);
    \node[note,below=0.18cm of g] {测试不通过\\则拒绝注册};
    \node[note,below=0.18cm of n] {能力固化};
  \end{tikzpicture}
  \end{center}

  \vspace{0.2em}
  \begin{columns}[T,onlytextwidth]
    \begin{column}{0.48\textwidth}
      \begin{block}{S1\quad 工具自扩展}
        \footnotesize
        注册接口强制要求提供测试文件；缺失即返回错误。
        \textbf{未经验证的自建工具会将其错误结果传播至后续全部运行，
        且难以归因，其危害大于缺少该工具。}
      \end{block}
    \end{column}
    \begin{column}{0.48\textwidth}
      \begin{block}{S2\quad 能力测量}
        \footnotesize
        16 题留出基准，全部答案来自本项目自身的完整枚举。
        主指标为 \textbf{通过题数 / 总题数}；
        \textbf{因工具缺失而跳过者计为未通过}。
      \end{block}
    \end{column}
  \end{columns}

  \vspace{0.3em}
  \begin{block}{S3\quad 策略先验}
    \footnotesize
    每次算子比较自动入库；后续以一次查询替代多轮测量。
    推荐须满足\textbf{跨 $\ge 2$ 次独立实验一致}且\textbf{均值差超过噪声地板}，
    否则返回 \code{inconclusive} 且不给出推荐值。推荐不跨 $n$ 迁移。
  \end{block}
\end{frame}

%% ---------------------------------------------------------------- 能力测量
\begin{frame}{能力增长的可测量证据}
  \begin{columns}[T,onlytextwidth]
    \begin{column}{0.46\textwidth}
      \centering
      \begin{tikzpicture}[x=0.97cm,y=0.044cm]
        %% 坐标轴
        \draw[black!45,-{Stealth[length=2mm]}] (0,0) -- (0,112);
        \draw[black!45] (0,0) -- (5.75,0);
        %% 网格与刻度
        \foreach \y in {25,50,75,100}{
          \draw[black!12] (0,\y) -- (5.7,\y);
          \node[left,font=\tiny,black!60] at (0,\y) {\y};
        }
        \node[font=\tiny,black!60,rotate=90,anchor=south] at (-0.75,56) {基准通过率(\%)};
        %% 柱
        \foreach \i/\v/\cov in {1/76.9/1, 2/100/2, 3/100/2, 4/100/3, 5/100/3}{
          \fill[accent!55] (\i-0.28,0) rectangle (\i+0.28,\v);
          \node[above,font=\tiny] at (\i,\v) {\v};
          \node[below,font=\tiny] at (\i,-2) {\i};
          \node[below,font=\tiny,black!55] at (\i,-9) {\cov};
        }
        \node[font=\tiny,black!60] at (3,-19) {运行序号};
        \node[font=\tiny,black!60] at (3,-26) {下行数字为工具覆盖种类数};
      \end{tikzpicture}
    \end{column}

    \begin{column}{0.50\textwidth}
      \begin{block}{读数}
        \footnotesize
        \begin{tabular}{@{}ccc@{}}
          \toprule
          运行 & 通过率 & 工具覆盖 \\
          \midrule
          1 & 76.9\% & 1 \\
          2 & 100.0\% & 2 \\
          4 & 100.0\% & 3 \\
          5 & 100.0\% & 3 \\
          \bottomrule
        \end{tabular}
        \vspace{0.4em}

        首末对比：\key{+23.1 个百分点}，工具覆盖 \key{+2 种}。
      \end{block}

      \begin{block}{指标设计的取舍}
        \footnotesize
        若以“通过 / 实际尝试”为主指标，则因工具缺失而跳过者不计入分母，
        能力下降无法体现。改用绝对指标后，
        \textbf{能力退化在曲线上可见}。
      \end{block}
    \end{column}
  \end{columns}
\end{frame}

%% ---------------------------------------------------------------- 数学产出
\begin{frame}{领域产出：判定性结论}
  \begin{center}
  \scriptsize
  \begin{tabular}{@{}p{0.52\textwidth}p{0.36\textwidth}@{}}
    \toprule
    \textbf{结论} & \textbf{类型与作用域} \\
    \midrule
    循环图族在 $n\in\{39,42,43,44,45,46\}$ 上无解 & 已证明（全空间枚举） \\
    循环图族在 $n\in\{30\text{--}38,40,41\}$ 上有解 & 已证明（全空间枚举） \\
    循环图族存在性非单调（$n=33$ 有解，$34,35$ 无解） & 已证明 \\
    $R(4,6)$ 循环族在 $n=36\text{--}41$ 全无解 & 已证明（子族排除） \\
    \textbf{44/45 阶非循环阿贝尔族各仅 1 个} & 已证明 \\
    \bottomrule
  \end{tabular}
  \end{center}

  \vspace{0.6em}
  \begin{block}{智能体在运行中给出的群结构收窄}
    \footnotesize
    因 $Z_{44}$ 同构于 $Z_{11}\times Z_4$、$Z_{45}$ 同构于 $Z_9\times Z_5$，二者均为循环群，
    而循环族已被排除，故：
    \begin{itemize}
      \item $n=44$：待枚举的非循环阿贝尔族仅 $Z_{11}\times Z_2\times Z_2$，
        空间 $2^{23}=8{,}388{,}608$；
      \item $n=45$：仅 $Z_5\times Z_3\times Z_3$，空间 $2^{22}=4{,}194{,}304$。
    \end{itemize}
    \textbf{待枚举范围由 4 个族缩减为 2 个。}该结论已独立复算确认。
  \end{block}

  \vspace{0.2em}
  \footnotesize\color{black!65}
  说明：$R(5,5)$ 的已知界未被改进。本项目的产出集中于智能体能力层面，
  数学结论为其副产物。
\end{frame}

%% ---------------------------------------------------------------- 演示
\begin{frame}[fragile]{现场演示}
  \begin{block}{运行入口}
    \begin{semiverbatim}
\small
cd D:\\myagent
.\\marsa.ps1 doctor     \ \ \ \# 六项环境自检
.\\marsa.ps1 tools      \ \ \ \# 已固化能力与能力曲线
.\\marsa.ps1 run "用 ramsey_compare 比较 n=44 上两个算子，
                 各 4 个 seed、8 秒预算"
    \end{semiverbatim}
  \end{block}

  \vspace{0.3em}
  \begin{block}{预期观察}
    \footnotesize
    \begin{enumerate}
      \item 智能体先调用 \code{ramsey\_prior} 查询先验，再决定是否重新测量；
      \item \code{ramsey\_compare} 返回 \code{inconclusive} 时，
        其\textbf{照录“无定论”}，不改写为“略优”。
    \end{enumerate}
    \vspace{0.3em}
    \textbf{该行为即为设计目标}：比较结果落在种子噪声内时，
    系统拒绝给出优劣判断。
  \end{block}

  \vspace{0.3em}
  \footnotesize\color{black!70}
  备用方案：\code{.\textbackslash marsa.ps1 benchmark} —— 纯确定性、不调用语言模型、
  30 秒内完成，不受网络与接口状态影响。
\end{frame}

%% ---------------------------------------------------------------- 局限
\begin{frame}{局限与后续工作}
  \begin{columns}[T,onlytextwidth]
    \begin{column}{0.48\textwidth}
      \begin{block}{当前局限}
        \scriptsize
        \begin{tabular}{@{}p{0.40\textwidth}p{0.46\textwidth}@{}}
          \toprule
          项 & 状态 \\
          \midrule
          $R(5,5)$ 的界 & 未改进（仍为 43--46） \\
          上界（证明型） & 不支持，仅有构造工具 \\
          巨大族枚举 & 受算力限制（$2^{23}$ 可行，$2^{43}$ 不可行） \\
          统计显著性 & 由工具层保证，智能体自身不具备 \\
          长时程无人值守 & 已验证 15 轮 \\
          \bottomrule
        \end{tabular}
      \end{block}
    \end{column}

    \begin{column}{0.48\textwidth}
      \begin{block}{后续工作（按优先级）}
        \scriptsize
        \begin{enumerate}
          \item \textbf{S4 技能自蒸馏}：由失败路径提炼规则并改写自身技能，
            须配置回归门禁。该层风险最高。
          \item \textbf{Tier 2 基准}：以完整自主运行度量决策能力
            （现有基准仅覆盖工具层）。
          \item 完成两个待枚举族：
            $Z_{11}\times Z_2\times Z_2$ @ $n=44$；
            $Z_5\times Z_3\times Z_3$ @ $n=45$。
          \item 跨会话状态延续。
        \end{enumerate}
      \end{block}
    \end{column}
  \end{columns}

  \vspace{0.4em}
  \begin{alertblock}{一条由实测确立的工作原则}
    \footnotesize
    在智能体自行归纳的六处错误中，五处源于\textbf{间接推理}
    （必要条件误用、公式记忆、计数转述），一处于报告转述；
    而经\textbf{枚举与直接验证}得到的四条结论均未出错。

    \vspace{0.2em}
    因此：对“证明性”论述一律以工具复核；对工具直接给出的判定可以采信。
    \textbf{验证器是唯一仲裁者。}
  \end{alertblock}
\end{frame}

%% ---------------------------------------------------------------- 结论
\begin{frame}{结论}
  \begin{block}{}
    \normalsize
    本文构建了一个具备能力自扩展机制的数学研究智能体，并给出其能力增长的
    可测量证据。主要工作包括：
    \begin{enumerate}
      \item 在开源框架之上补充\textbf{确定性验证层}，使计算性断言可被精确裁决，
        并将统计纪律内置于工具接口；
      \item 实现\textbf{生长回路}，使运行中产生的代码经验经校验后固化为
        后续运行可直接调用的能力，并以留出基准度量其增益；
      \item 以三次真实自主运行给出实证：系统可自行构造工具、自我校验、
        在证据不足时拒绝断言，并在获得反证时撤回结论。
    \end{enumerate}
  \end{block}

  \vspace{0.6em}
  \begin{block}{}
    \normalsize
    \centering
    传统方法改进的是\textbf{解}；本系统改进的是\textbf{自身}。\\[0.35em]
    能力曲线表明：在无人工干预的条件下，
    系统可以做到\textbf{下一次运行优于上一次}。
  \end{block}

  \vspace{0.4em}
  \footnotesize\color{black!60}
  代码与运行记录：\texttt{github.com/Bitwild4ever/marsa-ramsey-agent}
\end{frame}

\end{document}
